\section{Model-Based RL 的核心思想}

\subsection{什么是 Model-Based RL}

与 model-free 方法不同，model-based RL 显式地学习环境的\textbf{动力学模型}（dynamics model），然后利用该模型来做决策。

\begin{importantbox}{动力学模型}
学习一个函数 $f_\phi(\mathbf{s}, \mathbf{a})$ 来近似环境的状态转移 $p(\mathbf{s}' | \mathbf{s}, \mathbf{a})$：
$$f_\phi(\mathbf{s}, \mathbf{a}) \approx \mathbf{s}'$$
训练目标是最小化预测误差：
$$\min_\phi \sum_i \|f_\phi(\mathbf{s}_i, \mathbf{a}_i) - \mathbf{s}_i'\|^2$$
其中数据 $\mathcal{D} = \{(\mathbf{s}, \mathbf{a}, \mathbf{s}')_i\}$ 来自与环境的交互。
\end{importantbox}

\begin{figure}[H]
\centering
\includegraphics[width=0.9\textwidth]{slides-images/slide-05.jpg}
\caption{本讲大纲：Model-Based RL 的核心内容\protect\footnotemark}
\end{figure}
\footnotetext{Slides 第 5 页。}

\subsection{为什么要学习模型}

\begin{knowledgebox}{Model-Based vs Model-Free 的核心权衡}
\begin{itemize}
    \item \textbf{数据效率}：Model-based 方法通常比 model-free 方法数据效率高得多。学到的模型可以用来"想象"无数种可能的轨迹，而不需要真实执行。
    \item \textbf{代价}：模型不完美，基于不准确模型的决策可能导致性能下降。这就是\textbf{模型偏差}（model bias）问题。
\end{itemize}
\end{knowledgebox}

\subsection{本章小结}

Model-Based RL 的核心思路是"先学世界模型，再用模型做规划"。它在数据稀缺场景（如机器人）中特别有价值，但需要妥善处理模型误差。

